

{% Q3258082

\needspace{7\baselineskip}
\item \rtask \ponto{\pt} %
Julgue o próximo item.

Em Python, listas de elementos podem ser preenchidas por qualquer tipo de objeto, porém a quantidade de objetos que terão essas listas só poderá ser alterada durante a criação delas.

% F
{\setlength{\columnsep}{0pt}\renewcommand{\columnseprule}{0pt}
\begin{multicols}{2}
\begin{answerlist}[label={\texttt{\Alph*}.},leftmargin=*]
    \ti[V.]
    \ifnum\gabarito=1\doneitem[F.]\else\ti[F.]\fi % gabarito
\end{answerlist}
\end{multicols}
}

    % Alternativa correta: E - Errado
    %
    % Vamos entender o que está sendo abordado na questão. O tema central é a utilização de listas em Python, uma estrutura de dados fundamental em qualquer linguagem de programação moderna, especialmente relevante para o Cargo de Analista Judiciário - Tecnologia da Informação.
    %
    % No Python, listas são estruturas de dados dinâmicas que podem conter elementos de qualquer tipo. Isso significa que, diferentemente de arrays em outras linguagens de programação que têm tamanho fixo, as listas em Python podem ser modificadas a qualquer momento, permitindo a adição ou remoção de elementos mesmo após a sua criação.
    %
    % Vamos ao resumo teórico:
    %
    % Em Python, uma lista é definida entre colchetes, como em minha_lista = [1, 2, 3]. As operações mais comuns que podem ser realizadas em listas incluem:
    %
    %     Adicionar elementos: Usando métodos como append() para adicionar um único elemento no final da lista ou extend() para adicionar múltiplos elementos.
    %     Remover elementos: Utilizando remove() para remover um elemento específico ou pop() para remover pelo índice.
    %     Alterar elementos: Atribuindo um novo valor a uma posição específica, por exemplo, minha_lista[0] = 10.
    %
    % Com base nisso, a afirmação de que "a quantidade de objetos que terão essas listas só poderá ser alterada durante a criação delas" está incorreta. As listas em Python são dinâmicas, permitindo modificações de tamanho a qualquer momento durante a execução do programa.
    %
    % Agora, analisando a alternativa:
    %
    %     Alternativa E - Errado: Esta é a alternativa correta, pois a afirmação da questão está incorreta quanto à flexibilidade na manipulação das listas em Python, que são dinâmicas e permitem adições e remoções de elementos a qualquer momento.
    %
    % Espero que agora você tenha entendido por que a resposta correta é a alternativa E. As operações dinâmicas de listas são essenciais para qualquer desenvolvedor Python, especialmente para alguém no cargo de Analista Judiciário em Tecnologia da Informação.
}


